\chapter{Results and Discussion}
\label{chap:results_and_discussion}
This chapter presents the experimental setup and empirical evaluation of the methodology described in Chapter~\ref{chap:methodology}. Section~\ref{sec:results-overview} recaps the experimental design. Section~\ref{sec:results-headline} reports the headline per-condition metrics. Section~\ref{sec:results-stats} gives the omnibus and pairwise statistical analysis. Section~\ref{sec:results-findings} presents the three findings that follow from that analysis. Section~\ref{sec:results-failure} looks at per-bug-type failure patterns, and Section~\ref{sec:results-discoveries} records analysis-time discoveries that affect how the headline numbers must be read. Section~\ref{sec:results-decisions} explains the resulting change to which condition is designated the dissertation's proposed method. Section~\ref{sec:results-feedback} responds directly to the questions raised at the interim demonstration regarding automation, artificial data, and label leakage. Section~\ref{sec:results-limitations} states the study's limitations, and Section~\ref{sec:results-summary} closes the chapter.

\section{Experimental Overview}
\label{sec:results-overview}
Two experiments are reported. \textbf{EXP001}, a 16-clip pilot, compared uniform sampling with a naive prompt (A0B0) against pixel-only keyframe selection with the dual-path prompt (A1B1), and motivated a pre-registered power analysis that fixed the formal trial's scale (Section~\ref{sec:exp_design}). \textbf{EXP002}, the formal trial, is a $3\times2$ factorial ablation of frame selection (A0 uniform / A1 pixel-only / A3 pixel+DOM fused) against prompting (B0 naive / B1 dual-path). Four of its cells, A0B0, A1B1, A3B0, and A3B1, were run against all \textbf{120} clips of \texttt{dataset\_v3} (60 normal, 60 anomalous, balanced across six bug types), each clip repeated 3 times per condition, for \textbf{1{,}440} total judge calls. All calls used the single pinned judge model \texttt{gemini-3.1-flash-lite} at temperature \num{0.0}, and cost \textbf{\$1.71} in total, within \SI{0.1}{\percent} of the amount projected from a 24-call token-usage probe taken before the full run. This confirms that per-condition token cost scales predictably with frame count. All results below are for EXP002 unless stated otherwise; EXP001's pilot result is given in Table~\ref{tab:pilot-results} (Section~\ref{sec:exp_design}) and is not repeated here.

\section{Headline Results}
\label{sec:results-headline}
Table~\ref{tab:headline-results} reports accuracy (mean $\pm$ SD across the 3 repeats, with a bootstrap 95\% CI over the 120 clips), false-positive rate and anomaly recall (computed on the per-clip majority vote; see the aggregation-unit caveat in Section~\ref{sec:results-discoveries}), mean total tokens, and median latency, for each of the four conditions.

\begin{table}[htbp]
\centering
\caption[Headline per-condition results, formal trial ($n=120$)]{Headline results, formal trial ($n=120$ clips, 3 repeats per clip per condition). Accuracy is mean $\pm$ SD across repeats with a 10,000-resample bootstrap 95\% CI over clips (seed \texttt{20260728}); FPR and recall are on the per-clip majority vote. Latency is the median (see Section~\ref{sec:results-discoveries} for why the mean is unsafe to report for A3B1).}
\label{tab:headline-results}
\small
\begin{tabular}{@{}lccccc@{}}
\toprule
Condition & Accuracy $\pm$ SD [95\% CI] & FPR ($n{=}60$) & Recall ($n{=}60$) & Mean tokens & Median lat.\ (s) \\
\midrule
A0B0 (uniform + naive)             & 0.753 $\pm$ 0.005 [0.675, 0.828] & 0.167 & 0.667 & 6697 & 1.097 \\
A1B1 (pixel keyframes + dual-path) & \textbf{0.908} $\pm$ 0.000 [0.850, 0.958] & \textbf{0.000} & \textbf{0.817} & 4031 & 0.983 \\
A3B0 (fused keyframes + naive)     & 0.756 $\pm$ 0.005 [0.678, 0.831] & 0.117 & 0.633 & \textbf{3396} & 0.952 \\
A3B1 (fused keyframes + dual-path) & 0.842 $\pm$ 0.008 [0.775, 0.903] & 0.000 & 0.683 & 4072 & 0.903 \\
\bottomrule
\end{tabular}
\end{table}

A0B0's 6-frame, no-selection baseline reaches \SI{75.3}{\percent} accuracy with a false-positive rate of \SI{16.7}{\percent} (10 of 60 normal clips misclassified as anomalous). A1B1, the pilot's original ``ours'' condition, is the strongest of the four cells on every headline metric: \SI{90.8}{\percent} accuracy, zero false positives, and the highest anomaly recall at \SI{81.7}{\percent}, on 40\% fewer tokens than the baseline (4{,}031 vs.\ 6{,}697). A1B1 also reports the correct \texttt{defect\_type} on 31 of its 49 correctly identified anomalous clips.

\section{Statistical Analysis}
\label{sec:results-stats}

\subsection{Omnibus test}
An omnibus Cochran's Q test over the $120\times4$ per-clip correctness matrix rejects the null hypothesis that all four conditions perform identically: $Q=23.30$, $df=3$, $n=120$, $p=3.5\times10^{-5}$. This licenses the pairwise comparisons below.

\subsection{Pairwise accuracy comparisons}
Table~\ref{tab:mcnemar} reports exact McNemar tests on every pair of conditions, Holm-Bonferroni corrected across the six comparisons, with Cohen's $g$ as the accompanying effect size.

\begin{table}[htbp]
\centering
\caption[Pairwise McNemar tests, Holm-corrected]{Pairwise exact McNemar tests on the per-clip majority vote, Holm-Bonferroni corrected across 6 comparisons. $n_\text{disc}$ = number of discordant clips; the split gives how many of those favour the first- and second-named condition respectively.}
\label{tab:mcnemar}
\begin{tabular}{@{}lccccc@{}}
\toprule
Contrast & $n_\text{disc}$ & Split & $p$ & $p_\text{Holm}$ & Cohen's $g$ \\
\midrule
A0B0 $\to$ A1B1 & 21 & 1 / 20 & $2.1\times10^{-5}$ & $\mathbf{0.000126}$ & 0.45 \\
A0B0 $\to$ A3B0 & 23 & 11 / 12 & 1.0 & 1.0 & 0.02 \\
A0B0 $\to$ A3B1 & 25 & 7 / 18 & 0.043 & 0.130 & 0.22 \\
A1B1 $\to$ A3B0 & 24 & 21 / 3 & $2.8\times10^{-4}$ & $\mathbf{0.001385}$ & 0.38 \\
A1B1 $\to$ A3B1 & 8 & 8 / 0 & 0.0078 & $\mathbf{0.031}$ & 0.50 \\
A3B0 $\to$ A3B1 & 24 & 7 / 17 & 0.064 & 0.130 & 0.21 \\
\bottomrule
\end{tabular}
\end{table}

Only three of the six contrasts survive Holm correction: A0B0 vs.\ A1B1, A1B1 vs.\ A3B0, and A1B1 vs.\ A3B1. Every significant contrast involves A1B1. Neither A0B0 vs.\ A3B0 (fused selection's effect under a fixed naive prompt) nor A3B0 vs.\ A3B1 (the dual-path prompt's effect under fixed fused selection) reaches significance after correction.

\subsection{Efficiency comparisons}
Wilcoxon signed-rank tests (Holm-corrected) on per-clip total tokens, latency, and cost show that every selector-based condition (A1B1, A3B0, A3B1) beats the uniform baseline A0B0 on all three efficiency metrics at $p_\text{Holm} < 10^{-6}$ with rank-biserial correlation of \num{1.00}: every one of the 120 clips individually used fewer tokens under selection than under uniform sampling. Among the two dual-path conditions, A1B1 is both more accurate than A3B1 (Table~\ref{tab:mcnemar}) and cheaper on tokens by median (4{,}036 vs.\ 4{,}073), so A1B1's accuracy advantage is not bought with extra inference cost.

\section{Three Findings}
\label{sec:results-findings}
The statistical analysis above supports three findings, two of them negative for the pipeline's originally intended full design.

\paragraph{Finding 1: the fused (A3) selector buys no accuracy.} A0B0 $\to$ A3B0 moves accuracy from \SI{75.3}{\percent} to \SI{75.6}{\percent}, a change indistinguishable from noise ($p_\text{Holm}=1.0$, $g=0.02$). It does halve token usage (6{,}697 $\to$ 3{,}396, $-$\SI{49}{\percent}), but that saving comes from sending 3 frames instead of 6, something pixel-only selection (A1) also achieves, and cannot be attributed specifically to the DOM-fusion mechanism, since A0B0 vs.\ A3B0 is the contrast that isolates frame selection under a fixed prompt and shows no accuracy effect.

\paragraph{Finding 2: fusion actively harms the method once the dual-path prompt is in use.} A1B1 (\SI{90.8}{\percent}) outperforms A3B1 (\SI{84.2}{\percent}), $p_\text{Holm}=0.031$, $g=0.50$, and \textbf{all 8 discordant clips favour A1B1; none favour A3B1}. Holding the prompt fixed at B1, per-bug-type recall drops from pixel-only to fused selection on \texttt{flicker\_seizure} (8/10 $\to$ 5/10, the largest single-type change and the source of 3 of the 8 discordant clips) and \texttt{flyaway\_offset} (10/10 $\to$ 8/10), with \texttt{blur\_unreadable}, \texttt{skew\_distortion}, and \texttt{displaynone\_collapse} each losing exactly one clip; \texttt{vanish\_invisible} is unaffected at 10/10 in both.

\paragraph{Finding 3: the headline A0B0 $\to$ A1B1 gain is real but aliased.} The largest and most significant contrast in the study, A0B0 $\to$ A1B1 ($p_\text{Holm}=0.000126$, false positives 10/60 $\to$ 0/60), moves both factors at once: selection changes from uniform to pixel-only \emph{and} the prompt changes from naive to dual-path in the same comparison. The clean, single-factor prompt contrast under fixed selection, A3B0 $\to$ A3B1, does not survive correction ($p_\text{Holm}=0.130$), and because the frozen design contains no A1B0 cell (pixel-only selection with the naive prompt), the corresponding clean selection contrast under fixed prompt is not available for the A1 selector either. The headline gain therefore cannot be split into ``how much is the frame selection'' and ``how much is the prompt'' from the four cells that were run.

\section{Failure Analysis}
\label{sec:results-failure}
\texttt{displaynone\_collapse} is a systematic blind spot rather than an occasional miss: it is answered correctly in only \textbf{2 of 40} clip-condition pairs across the four headline conditions (0/10 under A0B0, 1/10 under A1B1, 1/10 under A3B0, 0/10 under A3B1). A collapsed todo card, where the input and footer remain but the list itself is hidden, looks visually similar to a legitimately empty todo list, and nearby typing-related DOM mutations from the interaction script can lead the dual-path prompt's decision rule (``content changes coinciding with DOM mutations are normal'') to read the collapse as an ordinary user edit rather than a rendering defect. The pilot dataset contained only a single clip of this type, which also failed; at $n=1$ this could not have been told apart from an incidental miss, and it took the formal trial's 10 clips per type to show that the failure is systematic rather than incidental. This is flagged as a concrete target for future refinement of the decision prompt or the structural-signal set (Section~\ref{chap:conclusion}), not a result this dissertation claims to have solved.

\section{Analysis-Time Discoveries}
\label{sec:results-discoveries}
Several properties of the scoring pipeline and the cached results came to light during analysis and materially affect how the headline numbers above should be read. Each is disclosed here rather than folded silently into the reported figures.

The one schema-noncompliant response is not a parsing failure. One response, for clip \texttt{bug\_flicker\_seizure\_5}, condition A3B1, repeat 2, returned well-formed JSON whose \texttt{reason} field correctly described the defect, but used the out-of-schema verdict string \texttt{``defect''} instead of \texttt{``anomaly''}; the scorer treats any verdict outside \{\texttt{normal}, \texttt{anomaly}\} as incorrect. This did not change the clip's scored outcome, since the other two repeats of the same clip both independently answered \texttt{``normal''} and the majority vote was already wrong, but it means detection accuracy and schema-compliance rate are different quantities that must be reported separately (Section~\ref{sec:benchmarking}), not read as if a schema violation were evidence of a parsing bug.

Two different aggregation units are in play, and results depend on which is used. \texttt{accuracy\_mean} in Table~\ref{tab:headline-results} is the mean of three per-repeat accuracies, while FPR, recall, and every statistical test in Section~\ref{sec:results-stats} instead operate on the per-clip \emph{majority vote} across the three repeats. The two units can disagree in scale: A0B0 $\to$ A3B0 reads as $+0.3$ percentage points on the accuracy-mean scale but $+1/120 = +0.8$ percentage points on the majority-vote scale used by the McNemar test in Table~\ref{tab:mcnemar}, and A0B0's FPR is $10/60=\SI{16.7}{\percent}$ at the clip (majority-vote) level against $29/180=\SI{16.1}{\percent}$ at the individual-response level. Every table in this chapter states which unit it uses.

Mean latency is unsafe to report for A3B1. One clip under A3B1 took \SI{39.97}{\second}, an outlier well outside every other observation, which lifts the condition's \emph{mean} latency to \SI{1.261}{\second}, nominally above A0B0's \SI{1.140}{\second} mean, while its \emph{median} is \SI{0.903}{\second}, and the rank-based Wilcoxon test in fact places A3B1 as faster than A0B0. Reporting the mean here would state the opposite conclusion from the one the statistical test supports, so Table~\ref{tab:headline-results} reports medians throughout.

The Finding-1 null result is a net null, not an absence of effect. A0B0 $\to$ A3B0 is accuracy-neutral overall, but it flips 23 of 120 individual clips, with substantial offsetting swings by bug type: \texttt{blur\_unreadable} recall falls from 10/10 to 6/10 while \texttt{vanish\_invisible} rises from 9/10 to 10/10, and \texttt{skew\_distortion} falls from 6/10 to 3/10 while \texttt{flyaway\_offset} rises from 6/10 to 8/10. ``Fused selection contributes nothing'' would misstate this; the accurate summary is that it has no net accuracy effect, at roughly half the token cost, with large offsetting per-type swings underneath.

The direction of Finding 2 contradicts the fusion rule's own design rationale. The fusion rule (Section~\ref{sec:system_architecture-fusion}) down-weights pixel transitions that coincide with DOM activity specifically to \emph{favour} CSS-only defects such as \texttt{flicker\_seizure} over ordinary user edits, yet \texttt{flicker\_seizure} is exactly the category where fused selection loses the most ground under the dual-path prompt. No verified mechanism for this reversal is established in this dissertation; a frame-level comparison of what A1 and A3 actually select on the 8 discordant clips is identified as a specific, well-scoped piece of future work (Section~\ref{chap:conclusion}) rather than something this analysis resolves.

There is also a disclosed deviation from the project's statistical-tooling specification. The supervisor-set statistics rule (Section~\ref{sec:benchmarking}) specifies \texttt{statsmodels} for Cochran's Q and the Holm correction; \texttt{statsmodels} was absent from the locked environment used for this trial, so Cochran's Q was implemented directly and SciPy supplied the Wilcoxon test and chi-square tail probability instead. The two exact-McNemar implementations that previously existed independently in the trial-scoring and power-analysis scripts were consolidated into one shared function before the formal trial, so the pre-registered power model in Section~\ref{sec:exp_design} is provably a model of the same test that produced Table~\ref{tab:mcnemar}, rather than a nominally matching but independently written duplicate.

Finally, there is a dispersion gap against the project's own reporting rule. The reporting rule in Section~\ref{sec:benchmarking} requires mean $\pm$ SD [95\% CI] ($n$) for every metric, but the scoring script currently emits a bootstrap SD and CI for accuracy only; FPR, recall, and defect-type match are reported in Table~\ref{tab:headline-results} and above as point estimates. The missing dispersion figures were computed separately from the cached raw responses for this write-up. Folding that computation into the scoring script itself, so the figures are traceable in the machine-readable summary rather than recomputed ad hoc, is recorded as an outstanding follow-up; no re-run of the trial is required, since scoring is a pure function of the already-cached responses.

\section{Reframing the Proposed Method}
\label{sec:results-decisions}
The frozen experimental configuration originally designated A3B1, fused selection with the dual-path prompt, as ``the full proposed method,'' with A1B1 recorded as merely the pilot's earlier, simpler condition. Findings 1 and 2 (Section~\ref{sec:results-findings}) do not support that designation: the fused selector adds no measurable accuracy under a fixed prompt, and actively reduces accuracy relative to pixel-only selection once combined with the dual-path prompt it was intended to pair with. This dissertation therefore reports \textbf{A1B1, pixel-only keyframe selection combined with the dual-path decision prompt, as the proposed method}, and treats the DOM-fusion selector (A3) as a design that was implemented, pre-registered, formally tested, and \emph{falsified} by the formal trial, rather than as a contribution to claim. This is a post-hoc change to which condition is presented as the headline result, made after the data were seen; it is recorded here explicitly, with its date and reasoning kept alongside the original frozen designation in the experiment's configuration record (Section~\ref{sec:exp_design}), rather than silently substituted, so the change and its evidential basis can be checked.

A related decision was to \emph{not} run the additional A1B0 cell (pixel-only selection with the naive prompt) that would fully de-alias Finding 3. Doing so would cost around \$0.32 and would answer a real open question, but it would be a design change made in direct response to already-observed results within what was declared a frozen experiment, and the freeze policy adopted for this project (Section~\ref{sec:exp_design}) treats that as requiring a new, separately pre-registered experiment rather than an amendment to EXP002. The pre-registered four-cell design is accordingly left intact, and the resulting aliasing, together with its concrete consequence that the evidence for Finding 2 (fusion harms accuracy) rests solely on the B1 layer and is uncorroborated at B0, is carried into the limitations below rather than patched over.

\section{Response to Committee Feedback}
\label{sec:results-feedback}
Three points raised in discussion of this project's interim demonstration are addressed directly here, in addition to being reflected throughout Chapters~\ref{chap:methodology} and~\ref{chap:results_and_discussion}.

\subsection{Which parts were done by hand and which were automated}
Section~\ref{sec:division_of_labour} gives the full account: once launched, dataset generation, both quality gates, all three frame selectors, prompt construction, judge calls, response caching, and scoring run without per-clip human input, and ground-truth labels are derived automatically from the injected bug key rather than hand-annotated. What was decided by hand is the pipeline's design: the bug taxonomy, the factorial ablation and which cells to run, the fusion rule and its parameters, the choice of statistical tests, the freeze policy, and, after the data were seen, the reframing decision in Section~\ref{sec:results-decisions} and the decision not to run A1B0. An AI coding assistant was used to implement modules against these hand-specified designs and to help reconcile reported figures against the cached raw results; it was not used to choose the experimental design or to decide which findings to report. A complete, dated log of prompts, the assistant's contribution at each stage, and what was independently verified is given in Appendix~\ref{app:genai-log}.

\subsection{How far the artificial data limits the study, and how far this is standard practice}
The dataset used in this dissertation is entirely synthetic: all 120 clips come from injecting one of six CSS/DOM-level defects, by script, into a single reference implementation of TodoMVC (Section~\ref{sec:data_prep-generation}). This is a genuine limitation on external validity, and Section~\ref{sec:results-limitations} treats it as one rather than downplaying it. At the same time, automated, synthetic fault injection of this kind is the dominant methodology in the closest related literature reviewed in Chapter~\ref{ch:litreview}, precisely because manually producing and labelling defective interaction recordings at any useful scale is prohibitively labour-intensive~\citep{guirobust}; the most directly comparable recent benchmark, WUICC-bench, is itself built by an automated, LLM-driven mutation pipeline rather than by collecting real-world defects~\citep{zhang2026beyond}. In that sense, the use of synthetic data here follows, rather than departs from, standard practice in this area.

Two differences from that comparator are worth stating precisely, because they pull in opposite directions. First, WUICC-bench's synthetic samples pass a \emph{human verification} step after automated generation, retaining \SI{78.09}{\percent} of generated samples with inter-annotator agreement $\kappa \approx 0.722$~\citep{zhang2026beyond}. The present study's labels are validated only by the two automated gates described in Section~\ref{sec:data_prep-qg} (a pixel-difference screenshot gate and a video-level final-frame gate), with no equivalent human-verification pass to establish that an independent rater would agree each injected change looks like a defect a developer would file. This is a weaker form of validation than the comparator's, and it is disclosed as such. Second, the present study's defects are injected into \emph{live, recorded interaction video} rather than mutating a static screenshot pair, and the video-level gate (Section~\ref{sec:data_prep-qg}) verifies the defect is provably present in the decoded video frames actually sent to the judge, which is a stronger and more direct guarantee for the specific claim this dissertation makes, that a temporal defect is detectable from video, than a screenshot-only check would give. Taken together, the dataset's automated construction fits accepted practice in this literature, but its ground truth has not been independently human-verified the way the closest comparable benchmark's has, and its scope is limited to one reference application and six defect categories chosen for illustrative coverage rather than for representativeness of real-world defect distributions. Generalisation beyond this dataset is accordingly not claimed, and is listed explicitly in Section~\ref{sec:results-limitations}.

\subsection{How the ground-truth label is kept out of the judge prompt}
The full, code-level argument is given in Section~\ref{sec:label_leak} and is repeated here only in summary, since that section is most directly responsive to the question as it was raised. The clip manifest carries five label-bearing fields (\texttt{ground\_truth}, \texttt{is\_anomaly}, \texttt{bug\_type}, \texttt{bug\_label}, \texttt{expectation}), and these are read only by the scoring script, \emph{after} a judge response has already been generated and cached; they are never referenced anywhere in the request-construction path (\texttt{analyze\_clip()} and \texttt{build\_condition\_payload()} in \texttt{pipeline/run\_trial.py}), which reads only the clip's video and DOM-trace file paths. Both frame selectors take only pixel and DOM data as arguments, by construction, as stated directly in \texttt{pipeline/keyframes.py}'s own module docstring; the dual-path prompt's ``expected appearance'' text is a single generic string identical across all 120 clips and all six bug types, not derived from the per-clip \texttt{expectation} field; and the structural signals passed to the model describe the selected pixel/DOM transition, not the injection. Because this is a property of the code that builds every request rather than a claim about the study's intentions, it can, in principle, be checked independently by inspecting the two named functions and the manifest schema directly, which is why the explanation is given with file and function names rather than only in prose.

\section{Limitations}
\label{sec:results-limitations}
The study's largest and most significant result, A0B0 vs.\ A1B1, changes both frame selection and prompting at once and cannot be decomposed with the four cells that were run (Finding 3, Section~\ref{sec:results-findings}, and Section~\ref{sec:results-decisions}); the A1B0 cell that would resolve this was deliberately not added post hoc. Relatedly, the finding that fusion harms accuracy rests entirely on the A1B1 vs.\ A3B1 contrast, both under B1: with no A1-selector, B0-prompt cell run, it is not established whether an equivalent effect would appear under the naive prompt. \texttt{displaynone\_collapse} remains a systematic blind spot, detected correctly in only 2 of 40 clip-condition pairs across all four headline conditions (Section~\ref{sec:results-failure}), so the method as evaluated should not be presented as reliable for this defect class.

All 1{,}440 calls used one pinned judge model, \texttt{gemini-3.1-flash-lite}. This was a deliberate choice to remove judge identity as a confound in the factorial ablation (Section~\ref{sec:env_setup-model}), but it means none of the findings in this chapter have been shown to generalise to a different VLM, including the open-weight model originally scoped for comparison. Likewise, as discussed in Section~\ref{sec:results-feedback}, all data come from one TodoMVC (React) instance and six illustrative defect categories, validated by automated gates rather than by independent human verification; generalisation to other UI frameworks, applications, or naturally occurring defects is not established by this study.

Two further gaps are disclosed rather than smoothed over. Cochran's Q was implemented directly rather than drawn from \texttt{statsmodels} as the project's own reporting rule specifies (Section~\ref{sec:results-discoveries}), and the scoring script currently reports bootstrap dispersion for accuracy only, not for FPR, recall, or defect-type match, although those figures were computed separately from the same cache for this chapter. Cost figures, finally, assume a stable, confirmed rate card for the single pinned model at the time of collection; they are not a claim about cost at other providers, other models, or future pricing.

\section{Summary}
\label{sec:results-summary}
The formal 1{,}440-call factorial trial supports pixel-only keyframe selection combined with the dual-path decision prompt (A1B1) as the strongest of the four evaluated configurations, reaching \SI{90.8}{\percent} accuracy with zero false positives at lower token cost than the uniform-sampling baseline, and this dissertation accordingly reframes A1B1, not the originally intended DOM-fused selector, as its proposed method. Two of the study's three central findings are negative results for the design as originally scoped: DOM-fused frame selection adds no accuracy under a fixed prompt and measurably harms it once combined with the dual-path prompt, and both are reported as such rather than reframed as successes. The headline comparison against the naive baseline is real and statistically robust but aliased across two design factors at once, one systematic defect category remains essentially undetected, and the evaluation is confined to a single judge model and a single, synthetically generated reference dataset. Chapter~\ref{chap:conclusion} draws these results together against the dissertation's original aims and sets out the specific, scoped pieces of future work, most immediately the frame-level A1-vs-A3 comparison and the A1B0 de-aliasing cell, that follow directly from the limitations stated above.
